% 原首行额外间距的回归复现：普通正文接 eqdenote，再接正文/标题。
% 功能矩阵见 eqdenote-boundary / eqdenote-pagination；原失败记录保持。
\documentclass[degree-level=doctor,fontset=fandol]{hithesis}
\begin{document}
\mainmatter
\noindent TableBefore 中文边界 Agjp\par
\begin{eqdenote}
\item $D_{\mathrm{p}}$ & TableOne 平均粒子直径（m）；
\item $\psi$ & TableTwo 孔隙体积占总体积的百分比；
\item $C_2$ & TableThree 特征渗透性（m$^2$）。
\end{eqdenote}
\noindent TableAfter 中文接排 Agjp\par
\section{TableHeading 中文标题}
\noindent TableEnd 中文结束 Agjp\par
\clearpage
\noindent PlainBefore 中文边界 Agjp\par
\noindent 式中 $D_{\mathrm{p}}$——PlainOne 平均粒子直径（m）；\par
\noindent $\psi$——PlainTwo 孔隙体积占总体积的百分比；\par
\noindent $C_2$——PlainThree 特征渗透性（m$^2$）。\par
\noindent PlainAfter 中文接排 Agjp\par
\section{PlainHeading 中文标题}
\noindent PlainEnd 中文结束 Agjp\par
\end{document}
